\chapter{Breaking the Barriers: A Historical Analysis of Pioneering Women in Mathematics}\label{chap:Historychap}

\section{Introduction: The Silent History of Mathematical Women}
The traditional narrative of the history of mathematics, which is deeply rooted in global tradition, has often been presented as a lineage of male genius, stretching from figures like Euclid and Archimedes to Newton and Gauss. This portrayal reflects centuries of systemic exclusion, where women were legally or socially denied access to formal education, scientific patronage, institutional appointments, and publication—the very avenues for recognition in mathematics (\cite{Alic86}). This marginalization not only resulted in a loss of potential talent but also created an incomplete and distorted historical record, obscuring the immense creativity and perseverance that women exhibited in engaging with the discipline.

A thorough historical analysis is undertaken to address this topic. The focus is on the lives, intellectual environments, and specific scientific contributions of pioneering women who, through remarkable perseverance, strategic navigation, or sheer intellectual vigor, succeeded in making foundational contributions to mathematics. Organized chronologically, the discussion traces the evolution of female involvement from the private philosophical schools of antiquity to the contemporary university system, while emphasizing the specific social, political, and institutional barriers encountered by each woman. By detailing their struggles, significant breakthroughs, and the enduring stories of their resistance, the aim is to fully integrate these figures into the primary historical narrative of mathematics, illustrating that their achievements often resulted from overcoming challenges largely absent for their male counterparts.

The emphasis will be on the unique challenges faced in various eras: the quest for philosophical credibility in antiquity; the use of pseudonyms and informal support during the Enlightenment; the urgent demand for formal qualifications and professional legitimacy in the 19th century; and the struggle against academic sexism and political oppression in the 20th century. By highlighting the context of their lives, we not only honor their brilliance but also deepen our understanding of the significant cultural influences that have shaped the field of mathematics.

\section{Antiquity}
The earliest documented contributions of women to mathematics did not arise from formal academies, which were absent in the modern sense, but rather from philosophical and family-based schools. This intellectual freedom frequently stemmed from exceptional privilege or close associations with prominent male scholars, which made the sustainability of such legacies quite precarious.

\subsection{Hypatia of Alexandria (c. 360-415 AD)}
Hypatia is widely acknowledged as the first female mathematician whose life and contributions are well-documented. She thrived in late Roman Alexandria, the final stronghold of classical Hellenistic learning, and was the daughter of Theon, a prominent mathematician and philosopher affiliated with the Museum (\cite{Waithe1987}).

\subsubsection{Education and Work} Hypatia benefited from a unique advantage: a thorough education under the direct guidance of her father. Her intellectual prowess encompassed mathematics, astronomy, and philosophy. Her primary focus was on the important tasks of preservation and teaching. She devoted herself to editing, clarifying, and expanding upon the foundational texts of Greek mathematics. Notably, she worked on Diophantus's Arithmetica (advanced algebra), and Apollonius's Conics (geometry). In an era marked by intellectual decline, her efforts were crucial in making these complex works accessible to the next generation of scholars (\cite{Snyder2000}).
\subsubsection{Authority and Influence}
Hypatia of Alexandria is an example of a unique intersection of intellectual prowess and public authority in an era characterized by gendered constraints on women's roles. Rather than introducing groundbreaking theorems, her significance is rooted in her formidable presence as a public intellectual. As the head of the Neoplatonic school, Hypatia offered instruction in mixed-gender and multi-religious classes, drawing students from across the Roman Empire. Moreover, her role as a political advisor to the Roman Prefect Orestes underscores her influence beyond academia. Historical accounts indicate that her ability to "appear in public before the chief magistrates" was an unprecedented feat for women of her time (\cite{Osen1974}). Hypatia's authority derived from her comprehensive knowledge and exceptional rhetorical abilities, marking her as a pivotal figure in both intellectual and civic spheres.
\subsubsection{Struggles and Tragedy}
Hypatia's intellectual prominence ultimately became her fatal flaw. As Alexandria was consumed by a power struggle between the secular authority (Orestes) and the religious authority (Bishop Cyril), Hypatia, a pagan philosopher advising the Prefect, was seen as a threat. In 415 AD, she was brutally attacked and killed by a Christian mob. This shocking event is not just a personal tragedy but is often cited as a symbolic end to the flourishing of classical, pagan scholarship in Alexandria, and a chilling testament to the extreme vulnerability of a woman who dared to occupy a position of public, intellectual power (\cite{Dzielska1995}). Her death sent a clear message: for a woman, public scholarship could be an act of dangerous transgression.
\section{The Age of Enlightenment and private Patronage (18th Century)}
Centuries passed before women gained public recognition in mathematics. The environment shifted from philosophical schools to the wealthy salons and courts of the Enlightenment, where high society and intellect intertwined.
    \subsection{Maria Gaetana Agnesi (1718-1799)}
    Maria Gaetana Agnesi is one clear example of succces and limitations of the privilege necessary for female scholarship in the 18th century.
        \subsubsection{Privilege}
        she was the eldest of 21 children from a prominent Milanese silk merchant family, and her figure exemplifies the intersection of privilege and industriousness in the development of intellectual capabilities. Recognized as a child prodigy, Agnesi demonstrated linguistic mastery by the age of nine, a testament to her early and rigorous education. Her father, keenly aware of her potential, facilitated her intellectual growth by organizing frequent intellectual conversations that provided her with a platform to engage in debates on philosophy and science (\cite{Mazzotti2007}). 
        \subsubsection{Work}
        Agnesi's motivation for her greatest work, \textit{Instituzioni analitiche ad uso della gioventù italiana} (1748), was notably practical; it was conceived as a comprehensive textbook intended to aid her younger brothers in their studies. This book explores how Agnesi's privileged upbringing intertwined with her strong work ethic to shape her contributions to mathematics and education in the 18th century. This two-volume treatise was a grounbreaking publication on differential and integral calculus, which sinthesized the work of Newton and Leibniz making some of the most advanced mathematics of the time more accessible. Because of its clarity and comprehensive purpose, it earned her wide fame and many translations across Europe.

        Another Maria Gaetana Agnesi's notable legacy is the cubic curve known as the "Witch of Agnesi," defined by the equation $x^2y = a^2(a - y)$. This curve exemplifies analytical geometry and is often featured in calculus textbooks. The term, however, stems from a linguistic misunderstanding; Agnesi referred to it as "la versiera," meaning "turning" in 18th-century Italian, which sounds like "avversiera," one of the words for "witch" in Italian. In 1801, John Colson's translation mistakenly rendered it as the "Witch of Agnesi," a memorable yet inaccurate name. This highlights the challenges of translating mathematical concepts across cultures and the humorous errors that can arise (\cite{Kennedy1980}).

        \subsubsection{The Symbolic Professorship}
        Her remarkable success led Pope Benedict XIV's to appoint her as the chair of mathematics and natural philosophy at the University of Bologna in 1750. This recognition made her the first woman to hold a university chair in mathematics, representing a monumental and symbolic victory against the institutional barriers of the time. However, she likely chose not to take up the position, devoting her life instead to charitable work and religious piety. This decision underscores how even significant achievements were often hindered by personal and social obligations that diverted women from pursuing long-term professional careers (\cite{Mazzotti2007}). The position was open, but the path for women to consistently hold such role  was not truly forged until the next century.

\vspace{1em}

The institutionalization of science during the 18th and 19th centuries —characterized by the rise of scientific academies, specialized journals, and prominent universities—sharply defined and intensified the barriers against women. They were generally prohibited from enrolling in universities, presenting papers, or participating in scientific society meetings. As a result, women's contributions often occurred in obscurity, necessitating ingenious means of circumvention.

\subsection{Émilie du Châtelet (1706-1749)}
Gabrielle Émilie Le Tonnelier de Breteuil, Marquise du Châtelet, was a pivotal figure in the French Enlightenment. Denied formal access to institutions, she used her wealth and aristocratic status to cultivate her own intellectual environment, particularly at her château, where she engaged with her renowned lover, Voltaire. Her struggle was to attain recognition as both a philosopher and scientist in a society that confined aristocratic women to predominantly social roles.

\subsubsection{Work}
Du Châtelet's most renowned mathematical work is her comprehensive translation of Isaac Newton's \textit{Philosophiæ Naturalis Principia Mathematica}. However, her contributions extended well beyond mere translation; they encompass a significant physical and philosophical analysis. In her commentary, she provided her own insights into Newton's theories, defending them against the Cartesian perspectives that were prevalent in France. Notably, her key contribution to physics was the clarification of the concept of kinetic energy, establishing it as proportional to the square of the velocity ($E\propto mv^2$), rather than simply to the velocity ($E\propto mv$) as had been previously suggested by others (\cite{Zinsser2006}).
\subsubsection{Impact}
Her translation of Newton, published posthumously after her death from complications related to childbirth, remains the definitive French version and solidified her status as a key figure in the transmission of modern physical mathematics in Europe. Her brief yet brilliant career tragically highlights how the biological realities of womanhood, coupled with a lack of systemic support, often resulted in the premature end of female intellectual pursuits.
\subsubsection{Act of Transgression}
Du Châtelet faced both institutional and social restrictions, barred from the University of Paris and excluded from the intellectual circles of the Enlightenment. A notable example is her need to dress in male attire to attend Café Gradot, where leading scholars debated mathematics and philosophy (\cite{Sperling2017}). These cafés were essential for idea exchange, yet women were strictly prohibited, forcing Du Châtelet to physically transgress gender norms for access. This disguise symbolizes the broader pseudonymity women employed, like Sophie Germain's alias "Monsieur Le Blanc," highlighting that the main obstacles to women's scientific contributions were not their intellect, but the gendered barriers of 18th-century Parisian intellectual life. Such actions illustrate the lengths women had to go to participate in discussions that advanced scientific thought.

\subsection{Sophie Germain (1776-1831)}
Germain's formative years were shaped by the tumult of the French Revolution, which, ironically, afforded her the solitude necessary for self-study. Confined indoors, she greedily explored the contents of her father's library, mastering subjects such as mathematics and Latin. Despite her parents' efforts to deter her—by confiscating her clothing and denying her heat and light—she persevered, studying by candlelight while wrapped in blankets (\cite{Perl1978}).
\subsubsection{The Ruse of "M. Le Blanc"}
Sophie Germain's use of the pseudonym "Monsieur Antoine-Auguste Le Blanc" was pivotal in her mathematical career, allowing her to gain recognition in a male-dominated field. Initially, she adopted this identity to submit papers and receive materials from the exclusive École Polytechnique.
Her early work on analysis was submitted to Joseph-Louis Lagrange, a leading figure at the École. Lagrange was so impressed with the rigor and originality that he requested a meeting with the aloof "Le Blanc." When he discovered that this brilliant correspondent was a self-taught young woman, he chose to support her rather than dismiss her as a fraud. This recognition transformed their relationship into one of personal patronage, providing Germain with invaluable feedback and direction. However, this support came at a cost: her academic trajectory depended entirely on his goodwill, lacking institutional equality.

The second and more dramatic revelation occurred years later with Carl Friedrich Gauss, known as the "Prince of Mathematicians." Their correspondence began after Germain, using the pseudonym Le Blanc, submitted her research on number theory to him. 

The ruse was revealed during the Napoleonic Wars when Germain, concerned for Gauss's safety during the French siege of Jena in 1806, contacted General Pernety to ensure his protection. When the general mentioned that he was acting on behalf of "Mlle. Le Blanc," Gauss was confused and demanded clarification. Germain then revealed her true identity, explaining her disguise was necessary due to "the prejudices that pursue the sciences with so much harshness when a woman presents herself."

Gauss's response on July 30, 1807, remains a powerful testament to her abilities (\cite{gray1978}). He wrote:
\textit{"How can I describe my astonishment and admiration at seeing my esteemed correspondent Monsieur Le Blanc change into this celebrated person whom he presents as a brilliant example of what I would find difficult to believe. The love for the abstract sciences, in general, and for the mysteries of numbers in particular, is very rare... But when a person of the sex which, by our customs and prejudices, must encounter infinitely more obstacles than men to familiarize herself with these thorny researches, succeeds nevertheless in surmounting these barriers and penetrating the most important of them, she must unquestionably have the noblest courage, extraordinary talent, and superior genius."}

This emotional validation from a leading mathematician confirmed that her work was exceptional and worthy of praise. It highlights how the greatest minds recognized that genius transcended the barriers of sex.
\subsubsection{Work}

Despite her significant contributions to number theory, particularly in developing the theory of Sophie Germain primes related to Fermat's Last Theorem, her most notable external recognition stemmed from her research in elasticity theory. After two unsuccessful attempts, she received the prestigious prize from the French Academy of Sciences in 1816 for her paper on the vibration of elastic surfaces, making her the only woman to achieve this honor on her own merits. However, she was never awarded a seat at the Academy or granted an honorary degree, and she often had to depend on male colleagues to present her work. Her life serves as a poignant example of the struggle of a genius who lacked institutional support.